\documentclass[preview]{standalone}

\usepackage{amsmath}
\usepackage{amssymb}
\usepackage{bettelini}
\usepackage{stellar}
\usepackage{definitions}

\begin{document}

\id{measuretheory-integral-functions}
\genpage

\section{Integral Functions}

\includesnpt{fundamental-theorem-of-calculus}

\begin{snippet}{integral-functions-riemann-comparison}
    If \(f\) is Riemann integrable on a bounded interval \([a,b]\), its
    integral function
    \[
        F(x)=\integral[a][x][f(t)][t]
    \]
    is Lipschitz continuous because
    \[
        |F(x)-F(y)|
        =\left|\integral[x][y][f(t)][t]\right|
        \leq\sup_{t\in[a,b]}|f(t)|\,|x-y|.
    \]
    In particular, it is uniformly continuous.

    For an improper Riemann integral the integrand need not be bounded, so
    Lipschitz continuity may be lost. Continuity at an improper endpoint is
    instead built into the existence of the improper limit. For example,
    \[
        \integral[a][b][f(t)][t]
        =\lim_{x\to a^+}\integral[x][b][f(t)][t].
    \]
\end{snippet}

\begin{snippetproposition}{lebesgue-integral-function-continuity}{Continuity of a Lebesgue integral function}
    Let \(I\subseteq\realnumbers\) be an interval,
    \(f\in\mathcal L^1(I)\), and \(a\in\overline I\), where an endpoint may
    equal \(+\infty\) or \(-\infty\). Then the oriented integral function
    \[
        F(x)=\int_a^x f(t)\,dt
    \]
    is continuous on \(\overline I\). If \(I\) is bounded, then \(F\) is
    uniformly continuous.
\end{snippetproposition}

\begin{snippetproof}{lebesgue-integral-function-continuity-proof}{lebesgue-integral-function-continuity}{Continuity of a Lebesgue integral function}
    Let \(x_n\to c\in\overline I\). Up to the harmless orientation convention,
    write
    \[
        F(x_n)=\int_I f(t)\chi_{(a,x_n)}(t)\,dt.
    \]
    For every \(n\),
    \[
        |f\chi_{(a,x_n)}|\leq|f|\in\mathcal L^1(I),
    \]
    and the indicators converge pointwise to \(\chi_{(a,c)}\), except
    possibly at the single point \(c\). Dominated convergence gives
    \(F(x_n)\to F(c)\). Thus \(F\) is continuous. If \(I\) is bounded,
    \(\overline I\) is compact, so continuity implies uniform continuity.
\end{snippetproof}

\begin{snippetexercise}{integrability-on-closure-of-open-set-exercise}{Integrability on the closure}
    If \(A\subseteq\realnumbers\) is open and \(f\in\mathcal L^1(A)\), under
    which interpretation or additional hypotheses can one conclude that
    \(f\in\mathcal L^1(\overline A)\)?
\end{snippetexercise}

\begin{snippettheorem}{lebesgue-integral-function-almost-everywhere-differentiability}{Almost-everywhere differentiability of integral functions}
    Let \(I\subseteq\realnumbers\) be an interval and
    \(f\in\mathcal L^1_{\mathrm{loc}}(I)\). For \(a\in I\), define
    \[
        F(x)=\int_a^x f(t)\,dt.
    \]
    Then \(F\) is differentiable almost everywhere and
    \[
        F'(x)=f(x)
    \]
    for almost every \(x\in I\).
\end{snippettheorem}

\begin{snippetdefinition}{locally-integrable-function-definition}{Locally integrable function}
    Let \(A\subseteq\realnumbers^n\) be open. A measurable function \(f\) is
    \emph{locally integrable} on \(A\), written
    \(f\in\mathcal L^1_{\mathrm{loc}}(A)\), if every point \(p\in A\) has a
    neighborhood \(U\subseteq A\) such that
    \(f\in\mathcal L^1(U)\).
\end{snippetdefinition}

\begin{snippetexample}{reciprocal-function-locally-integrable-example}{A locally integrable function which is not integrable}
    On \(I=(0,+\infty)\), the function \(f(t)=1/t\) is locally integrable but
    does not belong to \(\mathcal L^1(I)\). For every fixed \(a>0\),
    \[
        F(x)=\integral[a][x][\frac1t][t]
    \]
    is nevertheless defined for every \(x>0\). The base point cannot be moved
    to either endpoint: both \(\int_0^xdt/t\) and
    \(\int_x^{+\infty}dt/t\) diverge.
\end{snippetexample}

\begin{snippetexercise}{local-integrability-compact-subsets-exercise}{Compact-set criterion for local integrability}
    Prove that a measurable function on an open set
    \(A\subseteq\realnumbers^n\) is locally integrable if and only if it is
    integrable on every compact set \(K\subseteq A\).
\end{snippetexercise}

\begin{snippet}{almost-everywhere-defined-functions-and-integrals}
    On a complete measure space, a function used in integration need only be
    defined almost everywhere: it can be extended arbitrarily on the
    remaining null set without changing its integral. Similarly, a function
    is essentially bounded when it is bounded outside a null set.
\end{snippet}

\begin{snippetexample}{almost-everywhere-bounded-unbounded-function}{A function bounded almost everywhere}
    Let \(\{q_n\}_{n\geq1}\) enumerate the rationals and define
    \[
        f(x)=
        \begin{cases}
            n & x=q_n,\\
            0 & x\notin\rationalnumbers.
        \end{cases}
    \]
    The function is unbounded on every interval but is bounded almost
    everywhere, since it vanishes outside the null set \(\rationalnumbers\).
\end{snippetexample}

\begin{snippetdefinition}{essential-supremum-definition}{Essential supremum}
    Let \(f\colon X\fromto[-\infty,+\infty]\) be measurable on a measure
    space \((X,\mathcal M,\mu)\). Its \emph{essential supremum} is
    \[
        \operatorname*{ess\,sup}_{X}f
        \triangleq
        \inf\left\{
            k\in[-\infty,+\infty]
            \suchthat f\leq k\text{ almost everywhere}
        \right\}.
    \]
\end{snippetdefinition}

\begin{snippet}{essential-supremum-interpretation}
    The essential supremum is the least upper bound after values on a null
    set are ignored. The infimum in its definition is attained.
\end{snippet}

\begin{snippetexercise}{essential-supremum-is-essential-upper-bound-exercise}{The essential supremum is attained}
    Let \(M=\operatorname*{ess\,sup}f\). Prove that
    \(f\leq M\) almost everywhere.
\end{snippetexercise}

\section{Exercise}

\begin{snippetexercise}{logarithmic-function-series-integrability-exercise}{Integrability of a logarithmic function series}
    Determine whether
    \[
        f(x)=\sum_{n=1}^{\infty}\frac{\ln x}{n(nx^2+1)}
    \]
    belongs to \(\mathcal L^1((0,+\infty))\).
\end{snippetexercise}

\begin{snippetsolution}{logarithmic-function-series-integrability-solution}{logarithmic-function-series-integrability-exercise}{Integrability of a logarithmic function series}
    The series converges pointwise on \((0,+\infty)\). It cannot converge
    uniformly on the entire interval because each term is unbounded near
    zero. It does, however, converge locally uniformly. Indeed, on
    \([a,+\infty)\), with \(a>0\),
    \[
        \left|\frac{\ln x}{n(nx^2+1)}\right|
        \leq\frac{|\ln x|}{n^2x^2}
        \leq\frac{K_a}{n^2}.
    \]
    Thus \(f\) is continuous, and in particular measurable.

    On \((1,+\infty)\) all terms are nonnegative and
    \[
        \sum_{n=1}^{\infty}\int_1^{+\infty}
        \left|\frac{\ln x}{n(nx^2+1)}\right|\,dx
        \leq
        \left(\int_1^{+\infty}\frac{\ln x}{x^2}\,dx\right)
        \sum_{n=1}^{\infty}\frac1{n^2}<+\infty.
    \]

    On \((0,1)\) the terms are nonpositive. Integration by parts gives
    \begin{align*}
        \int_0^1\frac{-\ln x}{n(nx^2+1)}\,dx
        &=\frac1{n\sqrt n}
        \left(
            \left[-\ln x\,\arctan(\sqrt n x)\right]_0^1
            +\int_0^1\frac{\arctan(\sqrt n x)}{x}\,dx
        \right)\\
        &=\frac1{n\sqrt n}
        \int_0^1\frac{\arctan(\sqrt n x)}{x}\,dx.
    \end{align*}
    The boundary term vanishes. Split the last integral at
    \(x=1/\sqrt n\). Using
    \(\arctan(\sqrt n x)\leq\sqrt n x\) on the first part and
    \(\arctan(\sqrt n x)\leq\pi/2\) on the second gives
    \begin{align*}
        \int_0^1\frac{-\ln x}{n(nx^2+1)}\,dx
        &\leq\frac1{n\sqrt n}
        \left[
            \int_0^{1/\sqrt n}\sqrt n\,dx
            +\int_{1/\sqrt n}^{1}\frac{\pi}{2x}\,dx
        \right]\\
        &=\frac1{n\sqrt n}
        \left(1+\frac\pi2\ln\sqrt n\right).
    \end{align*}
    The series of these bounds converges. Therefore
    \(f\in\mathcal L^1((0,+\infty))\). The splitting point
    \(1/\sqrt n\) preserves the summable factor \(n^{-3/2}\).
\end{snippetsolution}

\end{document}
